\documentclass{article}
\usepackage[utf8]{inputenc}
\setcounter{secnumdepth}{0} % Set the numbering depth to 0

\title{Autonomous Vehicle Training in CARLA}

\author{Alex Johnson, Julian Rodriguez}

\date{\today}

\begin{document}

\maketitle

\section{Introduction}
\paragraph{}
Below is a list of the processes used and lessons learned 
as we trained our autonomous vehicle in CARLA. After explaining
our choices and exactly what went right and wrong, we will 
provide general recommendations for autonomous training based
on our experiences

\section{What Went Wrong}

\subsection{Custom PPO vs SB3}
\paragraph{}

\subsection{Driving Rules}
\paragraph{}
Our final car appeared to behave erratically at first, as though it wasn't
properly obeying driving rules like staying in the speed limit or using the
right lanes. However, we realized this was because of the way CARLA was initialized.

The speed limits in many of the CARLA maps change very abruptly, and often to unrealistic values.
One example is in the city map, where the speed limit around a turn is 30 mph, but it quickly changes
to 90 mph during the next straight segment. This makes the car appear to lurch forward. However, this
is actually good behavior of the model, as it is trying to stay at near 95\% of the speed limit.

Something similar happened with waypoints. The car will follow waypoints well, but if these
points are not placed right, the car will seem break traffic rules. Overall, it is super important
for later validation that the original simulation is made as accurate as possible, even details that
might appear to be right on first glance.

\subsection{?}
\subsection{?}
\subsection{?}

\section{What Went Right}

\subsection{
WandB
}
\paragraph{}{}
Weights and Biases (WandB) is a platform for storing model metrics during
machine learning training. It uses a simple Python API that can be inserted
into a training script. It is free to use and helps us track the model's
losses, collision rate, episode length, and style adherence, create
visuals, and monitor our runtime as we train, without even rendering
the car. It provided a quantitative way to track training, with live
updates and remote access via the website. 
\subsection{Style Vectors}
\paragraph{}
The style vector idea performed very strongly when rendered. Small changes 
to the vectors, which are constrained to the interval [0, 1], resulted in large 
output fluxuations. We found that keeping speed minimal and safety maximal led
to the smoothest car. In the future, training might want to reduce the loss 
related to the style vectors to make the model fluctuate less, ie. make the 
style changes less extreme.
\subsection{PPO Alerts}
\paragraph{}
We initially wanted a small, parameterless model for the adaptive alert system. So, we used Gaussion Kernel Estimation to essentially
weight the alert parameters by scores the system has achieved with them in the past. Then, the model will randomly sample near points
with higher scores to get the alert values (with some room to explore). 

Unfortunately, this converged much slower and with lower accuracy
than having a gradient-based model. We chose PPO for our gradient model using the route difference between the AV and hero car as the loss, 
using the typical PPO clipped loss. This converged well during simulation of multiple drivers.

\subsection{?}
\subsection{?}

\section{Conclusion}

1 paragraph

\end{document}
